\chapter{Conclusion}
\label{ch:conclusion}

\section{Research Question Responses}
\label{sec:rq-responses}

\todo{Salvage the "Research Question Responses" section from
\texttt{overleaf\_latex/chapters/chapter8-Discussion.tex}. Answer each RQ from
\cref{sec:research-questions} directly and in order. RQ3 draws on
\cref{sec:tradeoff}.}

\subsection{RQ1 --- Sovereignty Framework Compliance}
\label{subsec:rq1}
\todo{Answer in two or three paragraphs, pointing to \cref{sec:study1}.}

\subsection{RQ2 --- Readiness Relative to AWS}
\label{subsec:rq2}
\todo{Answer, pointing to \cref{sec:study2}.}

\subsection{RQ3 --- The Sovereignty--Readiness Trade-off}
\label{subsec:rq3}
\todo{Answer, pointing to \cref{sec:tradeoff}.}

\section{Recommendations for EU Enterprises}
\label{sec:recommendations}

\todo{Salvage the four strategies from the old discussion chapter:
lift-and-shift to a high-readiness EU provider; sovereignty-first selection
for sensitive workloads; hybrid architecture; monitor STACKIT for 2027--2028
re-evaluation. Add the build-versus-buy recommendation for Alps Alpine
specifically.}

\section{Vendor Lock-in Implications}
\label{sec:lockin}

\todo{Salvage from the old discussion chapter: OpenStack as an
interoperability anchor, proprietary managed services as the lock-in vector,
and the migration barrier away from AWS. Ties directly to the recommendations
above.}

\section{Limitations}
\label{sec:limitations}

\todo{Salvage the six subsections from the old discussion chapter:
self-reported evidence base, static snapshot of a moving market, AWS baseline
bias, service selection tailored to one enterprise, no SLA or
operational-quality dimension, exclusion of commercial terms. Add: the build
branch is assessed analytically and was never deployed.}

\section{Future Work}
\label{sec:future-work}

\todo{Deployment to the top-ranked provider belongs here — it is the obvious
next step and framing it as future work is exactly right now that it is out of
scope. Also: LLM self-hosting feasibility as a live trial, and re-running the
assessment once the IPCEI-CIS investments land.}
